\chapter{Where to go from here}
\label{ch:nextsteps}

\begin{aside}
You've reached the end of the runtime story. Where you go
next is up to your project. Here are some pointers.
\end{aside}

\section{Build something}

The fastest way to internalise the library is to build a
real app. Pick something small that interests you: a
todo, a clock, a chat client, a file browser.

\section{Read the source}

Chapter~\ref{ch:readingsource} pointed at the structure.
Read it for real: the headers are the documentation.

\section{Contribute}

Chapter~\ref{ch:contributing} explained how. Bug reports,
small fixes, big features --- all welcome.

\section{Teach}

If a chapter helped you, blog about it. If something
confused you, write a clearer version. The book benefits.

\section{Stay current}

The web moves fast; the terminal moves slowly but does
move. Keep an eye on terminal protocol developments,
new image protocols, BiDi progress.

\section{Goodbye}

You opened a book. You read about a tiny universe of
bytes and cells. You can now build interfaces that ship
in those bytes. That's a lot, in any size of book.

Thank you for reading.

\section{Recap}

\begin{recap}
\begin{itemize}[leftmargin=*]
  \item Build, read, contribute, teach.
  \item Stay current with protocol developments.
\end{itemize}
\end{recap}

\section*{Workshop}

\begin{enumerate}[leftmargin=*]
  \item Decide your next project.
  \item Sketch the screen.
  \item Open an editor.
\end{enumerate}
